%-------------------------
% Henry Speiser Resume - Hardware
%-------------------------
\documentclass[letterpaper,10pt]{article}

\usepackage{latexsym}
\usepackage[empty]{fullpage}
\usepackage{titlesec}
\usepackage{marvosym}
\usepackage[usenames,dvipsnames]{color}
\usepackage{verbatim}
\usepackage{enumitem}
\usepackage[hidelinks]{hyperref}
\usepackage{fancyhdr}
\usepackage[english]{babel}
\usepackage{tabularx}
\usepackage{geometry}

\geometry{letterpaper, top=0.3in, bottom=0.3in, left=0.4in, right=0.4in}

\pagestyle{fancy}
\fancyhf{}
\fancyfoot{}
\renewcommand{\headrulewidth}{0pt}
\renewcommand{\footrulewidth}{0pt}
\setlength{\footskip}{4pt}

\urlstyle{same}
\raggedbottom
\raggedright
\setlength{\tabcolsep}{0in}

% Section formatting
\titleformat{\section}{
  \vspace{-8pt}\scshape\raggedright\large
}{}{0em}{}[\color{black}\titlerule \vspace{-4pt}]

% Custom commands
\newcommand{\resumeItem}[1]{
  \item\small{#1}
}

\newcommand{\resumeSubheading}[4]{
  \vspace{-1pt}\item
    \begin{tabular*}{0.97\textwidth}[t]{l@{\extracolsep{\fill}}r}
      \textbf{#1} & #2 \\
      \textit{\small#3} & \textit{\small #4} \\
    \end{tabular*}\vspace{-6pt}
}

\newcommand{\resumeSubHeadingListStart}{\begin{itemize}[leftmargin=0.15in, label={}]}
\newcommand{\resumeSubHeadingListEnd}{\end{itemize}\vspace{-2pt}}
\newcommand{\resumeItemListStart}{\begin{itemize}[leftmargin=0.2in, itemsep=0pt, topsep=2pt, parsep=0pt, partopsep=0pt]}
\newcommand{\resumeItemListEnd}{\end{itemize}\vspace{-5pt}}

%-------------------------------------------
\begin{document}

%---------- HEADING ----------
\begin{center}
    \textbf{\Huge \scshape Henry Speiser} \\ \vspace{1pt}
    \small Berkeley, California $|$
    \href{mailto:hspeiser@berkeley.edu}{hspeiser@berkeley.edu} $|$
    \href{https://linkedin.com/in/henryspeiser}{\textcolor{blue}{linkedin.com/in/henryspeiser}} $|$
    \href{https://henryspeiser.com}{\textcolor{blue}{henryspeiser.com}} $|$
    (650) 781-2710
\end{center}

%---------- EDUCATION ----------
\section{Education}
\resumeSubHeadingListStart
  \resumeSubheading
    {University of California, Berkeley}{Expected Graduation May 2028}
    {B.S. Electrical Engineering and Computer Science $|$ \textit{College of Engineering}}{GPA 3.75/4.00}
\resumeSubHeadingListEnd

%---------- EXPERIENCE ----------
\section{Experience}
\resumeSubHeadingListStart

    \resumeSubheading
    {Aleph Surgical}{San Francisco, CA}
    {Member of Technical Staff Intern}{September 2026 -- Present}
    \resumeItemListStart
      \resumeItem{Working on autonomous suturing, surgical teleoperation, and robot learning.}
    \resumeItemListEnd

  \resumeSubheading
    {Origami Robotics}{San Francisco, CA}
    {Robotics Engineer Intern}{May 2026 -- August 2026}
    \resumeItemListStart
      \resumeItem{Started on hand retargeting and traced poor performance back to the hand's kinematics rather than the algorithm. Took over the Rev.~1.5 21-DoF redesign, changing joint axes, finger geometry, and actuator layouts to better reproduce human grasps.}
      \resumeItem{Existing hand benchmarks were easy to game, so I built grasp-based ones in MuJoCo, sampling millions of object/hand poses to test whether designs could actually grasp and hold objects, then using the results to search a 30-parameter design space.}
      \resumeItem{Designed, manufactured, and assembled the next-gen hand in CNC'd aluminum and 316L stainless, including actuator sizing and 48V three-phase axial-flux motor controllers with magnetic encoders and CAN small enough to fit inside the fingers.}
      \resumeItem{Built the synchronized teleoperation/data-collection stack around Manus gloves, Vive trackers, Raspberry Pi cameras, and robot joint states; also worked on fingertip tactile sensing.}
      \resumeItem{Worked on S1, Origami's data collection glove, whose 20 encoded finger joints map 1:1 onto the robot hand so demonstrations need no retargeting; it also records fingertip/palm pressure and video in sync. Early units were sold to OpenAI.}
    \resumeItemListEnd

  \resumeSubheading
    {Tesla}{Palo Alto, CA}
    {Low Voltage Electrical Engineer Intern}{September 2025 -- December 2025}
    \resumeItemListStart
      \resumeItem{Worked on Cybercab low-voltage systems including the windshield controller, TPMS, displays, Bluetooth, and GNSS, building fixtures and root-causing failures through thermal and vibration testing, including camera fogging in a room-sized thermal chamber.}
      \resumeItem{Every test had its own automation stack, so I built a modular framework where engineers could plug reusable hardware-control modules into one orchestrator. It cut lifecycle stress-testing time by 25\% and was adopted across the org.}
      \resumeItem{Built Python pipelines that aggregated thousands of hours of electrical logs in real time and pushed results directly into Jira, making failures visible while tests were still running.}
    \resumeItemListEnd

  \resumeSubheading
    {Somo AI}{San Francisco, CA / Mexico City, MX}
    {Hardware \& Firmware Engineer Intern}{April 2025 -- September 2025}
    \resumeItemListStart
      \resumeItem{Joined when Somo was about five people and had no hardware prototype, so I designed the first one from scratch. That board became the platform for figuring out the product's electronics, packaging, sizing, and field requirements.}
      \resumeItem{From there I designed 12+ boards across ESP32, STM32, Nordic, and Raspberry Pi, often shipping one or two revisions every two weeks; integrated mmWave, thermal, IMU, PIR, Ethernet, current sensing, CAN/RS-485, BLE, and 915~MHz radios.}
      \resumeItem{Wrote the RTOS firmware, tested the finished products, and designed enclosures; thousands of the resulting boards are now deployed in hotels around the world, including 1,000+ production units of my PIR sensor.}
    \resumeItemListEnd

  \resumeSubheading
    {Atila Biosystems}{Sunnyvale, CA}
    {Electrical \& Firmware Engineer Intern}{October 2022 -- December 2023}
    \resumeItemListStart
      \resumeItem{Came in with one goal, to recreate a \$600 PCR-testing device for roughly \$25 so it could go into homes and low-resource settings. Designed the PCB, power system, thermal hardware, and rugged solar/battery/USB-C enclosure around low-cost parts.}
      \resumeItem{Wrote embedded C for real-time PID thermal cycling; designed the heater electronics and state machines around the narrow temperature window needed to amplify biological samples without overheating them.}
    \resumeItemListEnd

\resumeSubHeadingListEnd

%---------- PROJECTS & LEADERSHIP ----------
\section{Projects \& Leadership}
\resumeSubHeadingListStart

  \resumeSubheading
    {FRC Team 971 (Spartan Robotics)}{Mountain View, CA}
    {Electrical \& Software Captain}{August 2021 -- June 2024}
    \resumeItemListStart
      \resumeItem{Hard impacts kept killing connectors and sensor wiring, so I made it my mission to eliminate the problem. I designed 11 PCBs for encoders, Hall sensors, beam breaks, and ToF, plus power, CAN, and vision interfaces for early-access Jetson Orins.}
      \resumeItem{Also wrote autonomous, superstructure, collision-avoidance, and auto-aim software with camera/IMU/encoder fusion; our 2022 autonomous ranked 2nd worldwide out of roughly 4,000 teams.}
    \resumeItemListEnd

  \resumeSubheading
    {Supersonic Rocket}{Los Altos, CA}
    {Lead Engineer}{June 2024 -- September 2024}
    \resumeItemListStart
      \resumeItem{A friend and I wanted to learn rocket propulsion from first principles, so we built nearly the entire vehicle ourselves, including APCP motor grains, thrust stands, a machined converging-diverging nozzle, avionics, wireless ignition, nose cone, and fins.}
      \resumeItem{Static-fired and modeled the motor before flying to Mach 1.2 and 3,250~ft, with measured motor and flight performance closely matching our simulations.}
    \resumeItemListEnd

\resumeSubHeadingListEnd

%---------- SKILLS & INTERESTS ----------
\section{Skills \& Interests}
\begin{itemize}[leftmargin=0.15in, label={}, itemsep=-2pt, topsep=2pt]
  \small{
    \item \begin{tabular*}{\linewidth}{@{}l@{\extracolsep{\fill}}l@{}}
      \textbf{EDA}\quad KiCad, Altium, LTspice, ANSYS & \textbf{Protocols}\quad CAN/RS-485, SPI/I2C, Sub-GHz (LoRa/OOK), BLE \\
      \textbf{CAD \& Sim}\quad Onshape, SolidWorks, CATIA, MuJoCo/URDF & \textbf{Firmware}\quad RTOS, Embedded Linux, Sensor Fusion, PID, FOC \\
      \textbf{Languages}\quad Python, C/C++, Assembly, Java & \textbf{Interests}\quad Ski Racing, Water Polo, Backpacking, Climbing, Movies \\
    \end{tabular*}
  }
\end{itemize}

%-------------------------------------------
\end{document}
